% strat_t2_usd — the unhedged US-dollar investor: FXGCF vs GCF timing.
% Numbers transcribed from tables/strat_t2_usd.pdf (generated by strategy.R).
\small
\begin{tabular}{l c c c c}
  \toprule
  Strategy & Mean (\%) & Vol (\%) & Sharpe & OOS $R^{2}$ \\
  \midrule
  $\FXGCF$-timed         & $2.68$ & $11.24$ & $0.24$ & $0.120$ \\
  $\GCF$-timed           & $1.95$ & $7.95$  & $0.25$ & $0.073$ \\
  Recursive-mean timing  & $1.60$ & $7.82$  & $0.20$ & --      \\
  Buy-and-hold           & $1.53$ & $7.72$  & $0.20$ & --      \\
  \bottomrule
\end{tabular}
\tabnotes{This table reports the performance of timing the unhedged US-dollar
  excess return \eqref{eq:rx-usd} of the 10Y global bond portfolio over
  2005--2024 ($n=221$ months). The $\FXGCF$-timed and $\GCF$-timed strategies
  use the recursive FX-adjusted and unadjusted global factors,
  $\FXGCF^{\mathrm{oos}}$ and $\GCF^{\mathrm{oos}}$, with the same
  mean--variance rule \eqref{eq:strat-weight} and equal-average-exposure
  reporting as the hedged strategy of \Cref{tab:strat-perf}. OOS $R^{2}$ is the
  \citet{campbellthompson2008} statistic of the corresponding factor forecast.
  The unhedged dollar bond premium is modest over the sample, so the dollar
  strategies improve on the passive portfolio but remain well below the
  currency-hedged investor.}
